\section{Experiment setup}


\subsection{Dataset}
We use \textsc{NormAd-eti} \cite{normad} for evaluation, a benchmark designed to assess the cultural adaptability of LLMs. The dataset contains 2.6K stories reflecting social and cultural norms from 75 countries, derived from the social-etiquette norms outlined in the Cultural Atlas.\footnote{\url{https://culturalatlas.sbs.com.au/}} Each story is associated with a country, a rule-of-thumb, and a ternary ground truth label $y \in$ \{Yes, No, Neither\} as shown in Figure \ref{fig:main_figure} (\textit{top}). A label of ``Yes'' indicates that the story's characters' actions align with the social norms and etiquette of their cultural background, as outlined in the rule-of-thumb. ``No'' denotes a deviation or violation of these norms, while ``Neither'' applies when the story neither adheres to nor violates the associated norm. Detailed statistics for each cultural bin categorized according to the Inglehart-Welzel cultural map are provided in Appendix Table \ref{tab:normad_stats}.



\subsection{Models}
We use the 7-9B variants of seven open-weight LLMs with varying levels of multilinguality. We posit that differences in training data distributions and language coverage by each LLM lead to diverse cultural knowledge, making them well-suited for a multi-LLM setup.\footnote{We use instruction fine-tuned LLMs, building on the findings of \textsc{NormAd-eti} \citep{normad}, which demonstrated that post-aligned models outperform their SFT couterparts in accuracy. We provide a detailed analysis in Appendix \ref{appendix:pre_post_alignment}.}

% \footnote{The original \textsc{NormAd-eti} paper \cite{normad} only considers English-centric LLMs.}.

\begin{itemize}[leftmargin=*, itemsep=2pt, parsep=-1pt]
 \item \textbf{English-centric:} \textsc{LLaMA-3} \cite{grattafiori2024llama3herdmodels}, \textsc{Gemma-2} \cite{gemma2}
 \item \textbf{Bilingual:} \textsc{EXAONE-3} \cite{exaone} (\textsc{en-ko}), \textsc{Yi-1.5} \cite{yi} (\textsc{en-zh}), \textsc{InternLM-2.5} \cite{internlm} (\textsc{en-zh})
 \item \textbf{Multilingual:} \textsc{Aya-23} \cite{aya23} (23 languages), \textsc{SeaLLM-3} \cite{seallm} (12 East Asian languages)
\end{itemize}
For multi-agent debate, we select \textsc{Gemma-2 27b} as our judge LLM accounting for its high single model accuracy compared to other LLMs and reasonable inference speed as detailed in Appendix \ref{appendix:judge_llm}. We set the default sampling temperature to 0.0, and employ 0.8 where multiple runs are required.\footnote{The HuggingFace model names are detailed in Appendix Table \ref{tab:huggingface_api}.}




\subsection{Evaluation Metrics}
\label{sec:evaluation}
Our primary evaluation metric is accuracy, calculated by comparing LLM responses to the ground truth labels.


% Given an LLM $\mathcal{M}$ and a dataset $\mathcal{D}$, we compare the ground truth label $y_i$ to the LLM response $\hat{y}^\mathcal{M}_i$ for accuracy scores. 

% \begin{equation}
%     \mathbf{Acc}(\mathcal{M}) = \frac{1}{|\mathcal{D}|}\sum_{\mathcal{D}} \mathbbm{1} [y_i = \hat{y}_i^\mathcal{M}]
% \end{equation}

% \input{table/decision_dynamics}


\definecolor{darkpurple}{RGB}{245, 66, 155}
\definecolor{bblue}{RGB}{66, 141, 245}
\definecolor{darkgreen}{RGB}{107, 201, 60}

\paragraph{Decision dynamics.}
For self-reflection and multi-agent debate, we analyze the dynamics of initial and final decisions made by LLMs. We aim to capture two key phases in this process: \textbf{\textcolor{darkpurple}{1) Initial Correctness}}: whether the LLM's initial decision is correct and \textbf{\textcolor{darkgreen}{2) Final Correctness}}: whether the LLM's final decision is correct. Specifically, for the multi-agent debate, we extend to capture additional key phase: \textbf{\textcolor{bblue}{3) Judge Correctness}}: whether the judge LLM's final decision is correct. If both agents already agree from \textbf{\textcolor{darkgreen}{2}}, this is determined by the correctness of the agreed-upon answer.
% We also consider both LLMs simultaneously.
% \textbf{\textcolor{bblue}{(2) Behavioral dynamic}:} whether the LLM changes its initial decision; and 
% We show the variations of all tested dynamics in Table \ref{tab:decision_dynamics}.



\paragraph{Cultural group parity.}
An essential aspect of cultural alignment is minimizing allocational harm and ensuring equitable performance across groups with diverse demographic attributes \cite{ramesh-etal-2023-fairness}. To operationalize this, we propose the concept of cultural group parity, which systematically evaluates how fairly our methods align with norms across different cultural groups \cite{wang-etal-2022-assessing}. Parity is defined as the state where methods exhibit comparable mean accuracies across cultural groups. Following \citet{petrov2023languagemodeltokenizersintroduce}, we set the baseline group $b$ as the group with the highest mean accuracy. For each cultural group $g$, we compute a parity premium relative to $b$. Formally:

\begin{equation}
    \mathbf{Parity}(g) = \frac{\mathrm{Acc}_g}{\mathrm{Acc}_b}
\end{equation}
A method achieves parity for $g$ with respect to $b$ if $\mathbf{Parity}(g)$ $\approx$ 1.
